%! root=./main.tex
\section{CFI Safety Analysis}

\heading{Setup.} We are given the following inputs:

\begin{itemize}[topsep=0pt,itemsep=-1ex,partopsep=1ex,parsep=1ex, leftmargin=*]
  \item $\pol : \mathbb{Z} \to \mathit{list}\ \mathbb{Z}$ --- the CFI policy
        % (maps each address to a list of valid target addresses),
  \item $\code : \mathit{list}\ \mathbb{Z}$ --- the original program code,
        % (a sequence of bytes),
  \item $bi \in \mathbb{Z}$ --- base index of the original code,
  \item $bi' \in \mathbb{Z}$ --- base index of the rewritten code,
  \item $tbi \in \mathbb{Z}$ --- policy table base index
        (the table maps each address to a list of permitted target addresses),
  \item $ai \in \mathbb{Z}$ --- abort index.
\end{itemize}

\noindent The CFI rewriter produces $(\codep,\, \tables)$ where:

\begin{itemize} [topsep=0pt,itemsep=-1ex,partopsep=1ex,parsep=1ex, leftmargin=*]
  \item $\codep : \mathit{list}\ (\mathit{list}\ \mathbb{Z})$ is the rewritten code, and
  \item $\tables : \mathit{list}\ (\mathit{list}\ \mathbb{Z})$ are the dispatch tables.
\end{itemize}

\subsection{Definitions}

\begin{definition}[Memory layout]
The three disjoint segments are:
\begin{align*}
  \OldText  &\;=\; \bigl[bi,\; bi + |\flatten(\code)|\bigr), \\
  \NewText  &\;=\; \bigl[bi',\; bi' + |\flatten(\codep)|\bigr), \\
  \DataSeg  &\;=\; \bigl[tbi,\; tbi + |\flatten(\tables)|\bigr).
\end{align*}
We require $ai \notin \NewText$ (the abort index is disjoint from the new
text segment).
\end{definition}

\begin{definition}[Block structure]
The rewriter partitions $\codep$ into blocks: each element $\codep[n]$
corresponds to exactly one instruction $\code[n]$.  We write $\block(n)$
for the $n$-th element of $\codep$, corresponding to the $n$-th
instruction of $\code$.
\end{definition}

\begin{definition}[Index translation --- old to new]
The function $\itoi : \mathbb{Z} \to \mathbb{Z}$ is defined by:
\[
  \itoi(i) \;=\;
  \begin{cases}
    \text{starting address of } \\ \block(i - bi) \text{ in } \NewText
      & \text{if } i \in \OldText, \\
    ai
      & \text{if } i \in \NewText, \\
    i
      & \text{otherwise.}
  \end{cases}
\]
Note that $\itoi$ maps only to block-start addresses or $ai$; it never
returns an index interior to a block.
\end{definition}

\begin{definition}[Index translation --- new to old]

The function $\itoio : \mathbb{Z} \to \mathbb{Z}$ is defined by:
\[
  \itoio(i') \;=\;
  \begin{cases}
    bi + n, \text{ where } n \text{ is the unique} \\
    \text{index such that }
    i' \in \block(n)
      & \text{if } i' \in \NewText, \\
    i'
      & \text{otherwise.}
  \end{cases}
\]
\end{definition}

\begin{definition}[Original policy]
$\Pol(i,\,j)$ holds if and only if $j \in \pol(i)$.
\end{definition}

\begin{definition}[Derived policy]
$\Polp(i',\,j')$ holds if and only if at least one of the following
conditions is satisfied:
\begin{enumerate}[label=(\arabic*), topsep=0pt,itemsep=-1ex,partopsep=1ex,parsep=1ex, leftmargin=*]
  \item $j' = ai$ \quad (transition to abort is always permitted),
  \item $\itoio(i') = \itoio(j')$ \quad (both addresses belong to the
        same original instruction's block), or
  \item $\Pol(\itoio(i'),\,\itoio(j'))$ \quad (the original policy
        permits the corresponding step).
\end{enumerate}
\end{definition}

\begin{definition}[Initial memory condition]
A memory state $m$ satisfies $\InitMem(m)$ if and only if:

\begin{itemize}[topsep=0pt,itemsep=-1ex,partopsep=1ex,parsep=1ex, leftmargin=*]
  \item $m$ contains $\flatten(\codep)$ starting at $bi'$,
  \item $m$ contains $\flatten(\tables)$ starting at $tbi$, with
        correct target endianness, and
  \item the regions $\NewText$ and $\DataSeg$ are nonwritable in $m$.
\end{itemize}
\end{definition}

\begin{definition}[Table well-formedness]
The tables are \emph{well-formed} if and only if every entry $e$ in
$\flatten(\tables)$ satisfies:
\[
  e = ai \qquad \text{or} \qquad
  \exists\, k \in \OldText.\ e = \itoi(k).
\]
\end{definition}

\begin{definition}[Valid trace]
A trace $t$ is \emph{valid} if and only if:
\begin{itemize}[topsep=0pt,itemsep=-1ex,partopsep=1ex,parsep=1ex, leftmargin=*]
  \item the entry-point of $t$ is $\itoi(k)$ for some $k \in \OldText$
        (i.e., a block-start address),
  \item the CPU begins in ARM mode (program counter divisible by $4$), and
  \item $\InitMem$ holds at the start of $t$.
\end{itemize}
\end{definition}

\begin{definition}[CFI-permitted step]
A single execution step from $i'$ to $j'$ within a trace is
\emph{CFI-permitted} if $\Polp(i',\,j')$ holds.
\end{definition}

\subsection{Lemmas}

\begin{lemma}[Endianness and mode stability]
Under $\InitMem$ and the assumption that no instruction in any block
modifies memory permissions or CPU endianness, the endianness of the CPU
and the ARM/Thumb mode (program counter alignment modulo $4$) are
preserved throughout any valid trace.
\end{lemma}

\begin{lemma}[Block-entry invariant]
If execution reaches an address $i' \in \NewText$, then $i'$ is a
block-start address; that is:
\[
  \exists\, k \in \OldText.\ i' = \itoi(k),
  \qquad \text{or} \qquad
  i' = ai.
\]
In particular, execution never enters a block at any address other than
its first instruction.

\medskip\noindent
\textit{Proof sketch.}
$\itoi$ maps only to block-start addresses.  All table entries are
block-start addresses or $ai$ by Definition~8.  Therefore no
control-flow transfer within the rewritten code can target a block
interior. \qed
\end{lemma}

\begin{lemma}[Block simulation --- case analysis]
Let $n$ be any index into $\code$, and suppose execution begins at
$\itoi(bi + n)$ with $\InitMem$ holding.  Then exactly one of the
following occurs:
\begin{enumerate}[label=(\arabic*),topsep=0pt,itemsep=-1ex,partopsep=1ex,parsep=1ex, leftmargin=*]
  \item \textbf{Unchanged instruction.}
        The $n$-th instruction was not altered; the block falls through
        to the next block, and the rewriter ensures that fallthrough is
        permitted by $\Pol$.

  \item \textbf{Changed, non-control-flow instruction.}
        The block either starts with a conditional branch directly to
        the next block (permitted by $\Pol$) or contains only
        fallthrough instructions.

  \item \textbf{Changed, static control flow.}
        The block contains a direct branch; the rewriter has verified
        that the target is permitted by $\Pol$ and has updated the
        branch target to account for inserted instructions.

  \item \textbf{Changed, dynamic control flow.}
        The block ends with a table-dispatch sequence: it loads a target
        address from the appropriate entry of $\tables$ and jumps to it.
        All table entries are $\itoi$-images of $\Pol$-permitted
        destinations, or $ai$.

  \item \textbf{Unsupported instruction.}
        The instruction is replaced by an unconditional branch to $ai$,
        which is permitted by clause~(1) of $\Polp$.
\end{enumerate}
In all cases, upon completion of the block:
\begin{itemize}
  \item the next program counter is either $ai$, a block-start address
        in $\NewText$, or an address outside $\NewText$,
  \item every step taken within the block is permitted by $\Polp$, and
  \item $\InitMem$ still holds (since no instruction modifies the
        nonwritable regions).
\end{itemize}
\end{lemma}

\subsection{Main Theorem}

\begin{theorem}[CFI Correctness]
For every valid trace $t$ (Definition~9), every execution step of $t$
is permitted by $\Polp$ until $t$ either exits $\NewText$ or raises a
hardware exception.  Moreover, $t$ can exit $\NewText$ only if the
original policy $\Pol$ permits the corresponding step.
\end{theorem}

\begin{proof}[Proof sketch]
By induction on the blocks of $t$.  The base case is the entry block,
which is a block-start address by Definition~9.

For the inductive step, Lemma~3 shows that each block executes only
$\Polp$-permitted steps internally and transfers control to either
$ai$, another block-start address, or an address outside $\NewText$.
Lemma~2 ensures that whenever control remains inside $\NewText$ it
arrives at a block-start, maintaining the induction hypothesis.
Lemma~1 ensures that ARM mode and endianness are preserved throughout,
so no implicit mode switch can bypass the analysis.  Finally,
$\InitMem$ is preserved by Lemma~3 (the nonwritable invariant), so
the hypothesis holds at each subsequent block.
\end{proof}

\heading{Remarks.}
The following points are left implicit above and may require additional
care in a proof assistant:

\begin{itemize}[topsep=0pt,itemsep=-1ex,partopsep=1ex,parsep=1ex, leftmargin=*]
  \item The well-definedness of $\itoio$ (every address in $\NewText$
        belongs to exactly one block) requires a disjointness lemma on
        the block partition.
  \item Lemma~3 is essentially a \emph{simulation relation} between old
        and new execution; it may be cleaner to state as a forward
        simulation depending on the chosen semantic framework.
  \item The ``hardware exception'' escape clause covers undefined
        behavior at the ISA level.  A separate assumption that no block
        contains invalid instruction encodings, or an additional terminal
        case, may be needed.
\end{itemize}